\chapter{Hardware, power and assembly}
\label{chap:hardware}

\section{Assembly inventory}

Before starting, record the information you can check for each component.
Do not infer a part number from a photograph or its similarity to another model.
If a detail such as a board revision or supply voltage cannot be confirmed from the
label, invoice or manufacturer's documentation, mark it as \emph{unverified}.

\begin{table}[H]
\centering
\caption{Basic assembly information.}
\label{tab:inventario-lector}
\begin{tabularx}{\textwidth}{@{}L{0.28\textwidth}L{0.31\textwidth}Y@{}}
\toprule
\textbf{Component} & \textbf{Information} & \textbf{Source} \\
\midrule
Host & Manufacturer, model, RAM and revision & Host, firmware or invoice \\
Operating system & Distribution, version and architecture & \code{/etc/os-release}, \code{uname -m} \\
Kernel & Full version & \code{uname -r} \\
Akida board & Model and revision & Label and BrainChip documentation \\
Adapter & Manufacturer, model and link type & Board markings and manufacturer's documentation \\
Cable & Type and orientation & Markings and adapter documentation \\
Host power supply & Voltage and power & Power supply label \\
Board or carrier power supply & Voltage, polarity and connector & Label and manufacturer's documentation \\
\bottomrule
\end{tabularx}
\end{table}

M.2, Mini PCIe, conventional PCIe and the Raspberry Pi 5 FFC connector are not
equivalent. Always check the connection type and cable orientation in the documentation
for the hardware you are using.

\section{Supported topologies}

\subsection{Linux host with a PCIe slot}

On a Linux PC, install the board in a compatible PCIe slot and power it as specified
in its documentation. Shut down and disconnect the host before inserting or removing
the board.

After assembly, continue with \cref{chap:host-driver}. If \code{lspci} does not detect
a new PCIe device, resolve that problem before configuring Python.

\subsection{Raspberry Pi 5 with an FFC-to-PCIe adapter}

Raspberry Pi 5 has a PCIe Gen~2~$\times$1 interface accessible through an FFC connector
\cite{rpi_pcie}. A conventional PCIe board needs an adapter that routes this link to
a compatible slot.

Even if the adapter has a physically longer slot, the Raspberry Pi 5 link remains
$\times$1.

\Cref{fig:montaje-aislado} shows the assembly used in the thesis.

\begin{figure}[H]
  \centering
  \includegraphics[width=0.92\textwidth]{hardware-carrier-detail-isolated.png}
  \caption[Raspberry Pi 5, carrier and AKD1000 assembly]{Assembly consisting of a Raspberry Pi 5,
  a PCIe adapter and the BrainChip AKD1000 board used in the thesis.}
  \label{fig:montaje-aislado}
\end{figure}

\section{Assembly}

If the manufacturer specifies a different connection sequence, follow that procedure.

\begin{enumerate}
  \item Shut down the host and disconnect the power supplies.
  \item Check the FFC cable orientation and the adapter or board supply voltage.
  \item Connect the FFC cable to the adapter in the specified orientation.
  \item Insert the AKD1000 fully into the carrier's PCIe slot.
  \item Inspect the connections before powering the system again.
\end{enumerate}

\begin{stopbox}{Electrical fault}
Disconnect power if you notice smoke, an unusual smell, excessive temperature or any
other sign of a possible power or assembly fault.
\end{stopbox}

\section{Waveshare carrier used in the thesis}

The thesis assembly used a Waveshare adapter connected to the Raspberry Pi 5 through
a 50~mm FFC cable. The board is marked \code{DC 12V} and
\emph{PCIe TO PCIe x1 Board (C)}.

Waveshare's current documentation describes this model as having a slot mechanically
compatible with PCIe $\times$1, $\times$4, $\times$8 and $\times$16 boards
\cite{waveshare_board_c}. The Raspberry Pi 5 link remains PCIe $\times$1 in every case.

The exact adapter item number was not retained. \Cref{tab:componentes-tfg} lists the
available identification details for the assembly.

\begin{figure}[H]
  \centering
  \includegraphics[width=0.78\textwidth]{hardware-powered-isolated.png}
  \caption[Powered assembly]{Assembly used during one of the thesis sessions.}
  \label{fig:montaje-alimentado}
\end{figure}

\section{Checks before power-on}

Before powering the assembly, check that the FFC cable has the correct orientation,
the board is fully seated and the supply voltages match those specified for the components.

\begin{stopbox}{Stop condition H0}
If you cannot confirm the voltage, polarity or orientation of a connection, do not
power the assembly until you have checked the relevant documentation.
\end{stopbox}
